\subsection{All six remaining two-by-three sources of order eight}

\begin{theorem}[Proved: the six order-eight two-by-three sources]
\label{CP2:A-eight-two-by-three-all-sources}
Let $M$ be any of the source tables with catalog index
\[
 492,\quad494,\quad495,\quad497,\quad499,\quad500.
\]
For every finite nonempty left-$R_3$, right-$M$ biset $S$,
all rational degree-two products of its collage vanish.
There is no state-count, boundedness, or Betti-number hypothesis.
\end{theorem}
\begin{proof}
All six monoids have identity $0$, band
$K=\{e_{ij}:i=0,1,\ j=0,1,2\}$, and one extra actor $a=7$.
The literal band labels are $e_{ij}=1+3i+j$.
Their complete tables are specified by the band product,
the identity, and the following data:
\[
\begin{array}{c|c|c|c}
\text{index}&\sigma_a&\tau_a&a^2\\ \hline
492&00&000&e_{00}\\
494&00&001&e_{00}\\
495&00&002&a\\
497&01&012&1\\
499&01&021&1\\
500&10&021&1
\end{array}                                                   \tag{1}
\]
Here $ae_{ij}=e_{\sigma_a(i),j}$ and
$e_{ij}a=e_{i,\tau_a(j)}$; a word lists a transformation's
values in increasing input order, and $1$ in the last
column denotes the global identity, whose literal label is $0$.

First consider the first five rows.  Let $P_c$ denote
the identity-adjoined $2\times c$ band.  The following
are actual surjective monoid maps, sending
$e_{ij}$ to $e_{i,\pi(j)}$:
\[
\begin{array}{c|c|c|c}
\text{index}&\text{target}&\pi&\rho(a)\\ \hline
492&P_3&012&e_{00}\\
494&P_2&001&e_{00}\\
495&P_2&010&e_{00}\\
497&P_3&012&1\\
499&P_2&011&1
\end{array}                                                   \tag{2}
\]
Each map takes the minimal left ideal $Me_{00}$
bijectively onto $P_c e_{00}$, with its actual left
$M$ action.  Consequently
Lemma~\ref{CP2:A-idempotent-quotient-bridge-connectivity}
applies: on the actual induced state quotient
$U=S\otimes_M P_c$, the bridge map is an isomorphism
on $H_1$ and a surjection on $H_2$.

We check the corner $H_2$ map explicitly, since three
of (2) collapse band columns.  Write
$v_1=j_1-j_0$, $v_2=j_2-j_0$ for the column
augmentation basis and $\chi$ for the row character
of $a$.  The balanced presentation of
Lemma~\ref{A:band-Tor} is
\[
 H_2(BM;\mathbb Q)=
 \mathbb Q\{v_1,v_2\}/\operatorname{im}(T-\chi I),            \tag{3}
\]
where $T$ is the column action of $a$.
In rows $492,497$ the relation matrix is zero and
the column map is the identity.
In row $494$ its image is $\mathbb Qv_1$,
and $\pi$ kills $v_1$ and preserves $v_2$.
In row $495$ its image is $\mathbb Qv_2$,
and $\pi$ kills $v_2$ and preserves $v_1$.
In row $499$ its image is $\mathbb Q(v_1-v_2)$,
and $\pi$ identifies $v_1$ and $v_2$.
Thus each actual corner map is an isomorphism on
second rational homology.  Its $H_1$ map is also
an isomorphism by the same connectivity lemma.

The natural two-block exact sequence now gives an
isomorphism on collage $H_2$: the bridge $H_2$
surjection identifies the preceding cokernels, while
the bridge and corner $H_1$ isomorphisms identify
the following kernels.  This is the same exact-sequence
argument as in
Proposition~\ref{CP2:A-band-quotient-collage-H2}, with
the corner check (3) supplied directly.
Its dual is the actual cohomological pullback.
Every $P_2/R_3$ target has zero degree-two products
by the bounded-six graph theorem, since its corner
projective dimensions give a degree-four homology bound.
Every $P_3/R_3$ target has zero degree-two products by
Theorem~\ref{CP2:A-two-by-three-band-right-zero-all-states}.
Their actual pullbacks prove the assertion for all
five rows, without assuming boundedness of the original.

It remains to treat row $500$.  Put $R=\mathbb QM$,
$B=\mathbb QP_3\subset R$, $e=e_{00}$,
$w=e_{10}-e_{00}$, and $q=(1-a)/2$.
There is the following complete projective left-$R$
resolution of the trivial module:
\[
 0\longrightarrow Re\xrightarrow{d_2}Rq
   \xrightarrow{d_1}Re\xrightarrow{\epsilon}\mathbb Q
   \longrightarrow0,                                        \tag{4}
\]
where $d_1$ is right multiplication by $w$ and $d_2$
is right multiplication by $e_{01}-e_{02}$.
Indeed $q^2=q$, $qw=w$, and
\[
 Re=\mathbb Q\{e_{00},e_{10}\},\qquad
 Rq=\mathbb Q\{q,e_{01}-e_{02},e_{11}-e_{12}\}.
\]
The map $d_1$ sends $q$ to $w$ and kills the two
band differences.  The injective map $d_2$ sends the
two displayed basis vectors of $Re$ to those differences.
These observations prove exactness and all maps are
left $R$ linear by their definition.

For $B$ use the pure-band resolution with degree-zero
term $Be$, degree-one term $Bq_0$, $q_0=1-e$, and
degree-two term $(Be)^2$.  Its degree-one differential
is right multiplication by $w$, and its degree-two
differential is right multiplication by
$e_{01}-e_{00}$ and $e_{02}-e_{00}$ on the two summands.
The inclusion $B\to R$ has a comparison map to (4)
which is the identity in degree zero, right multiplication
by $q$ in degree one, and
\[
 (u,v)\longmapsto (u-v)/2                                   \tag{5}
\]
in degree two.  To check the last assertion, use
$eq=0$ and
\[
 (e_{01}-e_{00})q=(e_{01}-e_{02})/2,\qquad
 (e_{02}-e_{00})q=-(e_{01}-e_{02})/2.
\]
Both resolutions augment to the identity on the
trivial module, so this chain comparison represents
the actual inclusion-induced map on Tor.

Take $L=\mathbb Q[S]$, retaining its commuting left
$R_3$ action.  There are no degree-three resolution
terms.  Thus the two second Tor modules are the
literal cycle kernels in degree two.  A target cycle
$t\in Le$ satisfies
\[
 t(e_{01}-e_{02})=0.
\]
Its pair $(t,-t)$ is a cycle in the pure-band complex,
since
\[
 t(e_{01}-e_{00})-t(e_{02}-e_{00})=0.
\]
Equation (5) sends that pair back to $t$.
This gives a natural left-$R_3$ linear section of
\[
 \operatorname{Tor}^{B}_2(L,\mathbb Q)
       \longrightarrow\operatorname{Tor}^{R}_2(L,\mathbb Q).
                                                               \tag{6}
\]

The trivial right $\mathbb QR_3$ module has projective
dimension one.  Since both source resolutions have
length two, the bridge spectral sequence in total
degree three has the single term
$\operatorname{Tor}^{\mathbb QR_3}_1
 (\mathbb Q,\operatorname{Tor}_2(L,\mathbb Q))$.
The split surjection (6) induces a surjection on
that term, hence on third bridge homology.
Both source corners have homology zero in degrees
three and four, and $BR_3$ is contractible.
The natural two-block sequence therefore identifies
that map with a surjection
\[
 H_4(B\mathcal C(R_3,S,P_3);\mathbb Q)
       \twoheadrightarrow
 H_4(B\mathcal C(R_3,S,M);\mathbb Q).
\]
The actual inclusion consequently induces an injection
on fourth cohomology.  Every product restricts to zero
on the pure-band collage by the all-state theorem,
so it was zero already.  This proves the final row.
\end{proof}

\begin{remark*}
The separate literal certificate binds all six full
tables to their catalog indices, all five quotient
maps and balanced-kernel calculations, and every
projective-basis and chain-comparison equation in
(4)--(6).  The proof is uniform in the biset; no
enumeration of state actions is a premise.
\end{remark*}
